\section{Clause Database Management}

Clause database management is another essential addition to a CDCL SAT solver. 
For larger and larger equations, and for equations with less valid assignents, 
we will tend to learn more clauses. While First UIPs can efficiently compute 
these learned clauses, with the added benefit of them being smaller than if we 
didn't use the first UIP condition, the number of learned clauses still has a 
memory impact. Although our test sets are of 3-SAT instances, so the initial 
clauses are small to begin with, and the learned clauses aren't as large as they 
could be, for larger SAT problems, the number of learned clauses can become exponential
in size. 

Although some people like to say that memory management isn't important because the 
hardware gets better and better, the day we have to stop worrying about the amount 
of memory being used is the day we have unlimited memory to work with. And until that 
day comes (which will be never), those people can shut up and go back to coding their 
2GB Chrome tab website. 
\footnote{Websites take up too much memory, much more than they need, 
and I will happily die on this hill.}

For this section, we will go back to using this paper:
\begin{center}
    \small{\url{https://www3.cs.stonybrook.edu/~cram/cse505/Fall20/Resources/cdcl.pdf}}
\end{center}
starting from section \textbf{4.5.4. Clause Deletion Strategies}.
The three (listed) methods described are:
\begin{itemize}
    \item \textit{n-order learning}: Only learning clauses with $n$ or fewer literals.
    \item \textit{m-size relevance-based learning}: Learning all clauses as long as they 
            currently imply assignments or are themselves a unit clause, only being
            only being discarded once they have $m$ or more Unassigned literals.
    \item \textit{k-bounded learning}: Clauses smaller than $k$ are kept permanently,
            but larger clauses are discarded once they stop being unit.
\end{itemize}
\noindent
The paper also mentions that \textit{k-bounded learning} and \textit{m-size relevance-based learning}
can be combined. You may wonder why we can't combine all three, however, notice that 
\textit{k-bounded learning} is an extension of \textit{n-order learning}, as they both keep 
clauses smaller than some threshold, but \textit{k-bounded learning} is more merciful to 
larger clauses.

The paper also describes a newer heuristic that is based on its \textit{activity} score. 
The heuristic is essentially ``VSIDS for clauses instead of variables.'' We will progress our 
way into implementing all of these algorithms, incrementally of course. Before we continue, it 
should be noted that modern SAT solvers don't have a universal clause deletion policy, and each 
one does its own thing, and once ours is mature enough, it will be no different.
\footnote{Does that make us all unique or all the same?}

\subsection{N-Order Learning}

\textit{n-order learning} is definitely the most aggressive clause learning/deletion policy of 
the bunch. I don't think I'd want to encounter it on the streets. It's very straightforward in 
idea: if a clause has more than $n$ literals, it won't get learned at all. This may raise some red 
flags, and I hope the one we're all thinking of is, how do we backjump? If we're not learning a clause, 
how do we know which decision level to backjump to? Well, just because we're not learning it doesn't 
mean we can't backjump to a higher decision level, it just means we won't add the new clause to our 
database. Does that put you at ease? Well it shouldn't because the reason why this was such a big 
red flag in the first place is that as of now, when we decide a new variable, we always start with 
True and rely on clause learning to force it False later. So when we backjump and start a new 
decision level, how do we start the new decision level as False if we didn't learn anything?

As of now, there is no way to change our course.
\footnote{This will be addressed in the next section.}
So what can we do? You're gonna find this funny (not really), we need to bring back our 
``tried second branch'' flag, or at least some form of it. Without it, and I cannot emphasize 
this enough, we will just be stuck in an infinite loop.

So how should we pick a value of $n$? There isn't a universal way to choose, so we will try 
the maximum clause length of the original CNF, and then try some numbers around it. Because 
our current testing sets are all instances of 3-SAT, it's not very fair, in my opinion, to 
test on different sizes of $n$ when the data is specifically tailored to be a fixed size. Because 
of this, we will go to our CNF database 
\footnote{\url{https://www.cs.ubc.ca/~hoos/SATLIB/benchm.html}}
and get a new testing set, ``Flat200-479'' which is a SAT encoding of a 3-color problem. There are 
600 variables and 2237 clauses, all satisfiable.
\footnote{As I was setting up the benchmarking test, my framework reported that the solver failed on 
this new test set. It turned out I gave it the wrong path, and after providing it the correct path, it 
passed, as we would expect. The shock of seeing the [FAIL] really threw me for a loop.}

Here are the baseline results (without n-order learning):
\begin{FVerbatim}
On all 1000 equations in uf20-91:

AMA:     274ms (~0.27 seconds).

On all 100 equations in Flat200-479:

AMA:     24,238ms (~24.24 seconds).

On all 1000 equations in UUF50.218.1000:

AMA:     1,629ms (~1.63 seconds).
\end{FVerbatim}
\noindent
We will now discuss our changes to our algorithm. Firstly, we need to compute the 
maximum clause length. As of now, we will say that the maximum clause length of the 
original equation will also be the maximum clause length of any learned clause, and then 
see how changing this number to lower or higher changes our runtime speed. To compute 
the maximum clause length, we will iterate through the original clauses and take the 
maximum size.

Next, we want to check how large the learned clause will be. We do this by checking 
the size of the \textit{seenIdx} array because it holds the variables that will 
become part of the learned clause. This avoids computing what the learned clause will 
be unnecessarily. For now, we will just backtrack one level in this case, then move onto 
backjumping.

We first need to remember to change our interface to use our old DecisionLevel struct:
\begin{lstlisting}
    public interface:

    private interface:
        DecisionLevel {
            int idx,
            bool triedOtherBranch
        }

        DecisionLevel[] trailStarts
        ...
\end{lstlisting}
\noindent
And of course we need to adjust our code to use the DecisionLevel struct instead of just an index.

The algorithmic changes are as follows:
\begin{lstlisting}
bool solve() {
    //setup
    ..

    while (True) {
        //assigning
        ...

        optional int[] contradiction = propagate(assignment)

        if (contradiction has a value) {
            numConflicts = numConflicts + 1

            //set up UIPs from conflict clause
            ...

            for (i = trail.size down to 0) {
                TrailEntry entry = trail[i]
                idx = abs(entry.lit)

                if (seen[idx]) {
                    //grab reason and update seen and levelCount
                    ...

                    delta[idx] = True
                    //unmark seen and update levelCount
                    ...
                }

                if (levelCount == 1) {
                    break
                }
            }

            if (seenIdx.size > maxSize) {
                while (trail.size > 0 && trail_starts.size > 0) {
                    DecisionLevel decision = trailStarts.last

                    while (trail.size > decision.idx + 1) {
                        int lit = trail.last.lit
                        assignment[abs(lit)] = Unassigned
                        var_to_trail[abs(lit)] = -1
                        erase(trail, trail.last)
                    }

                    if (decision.triedFalse == False) {
                        TrailEntry entry = trail.last

                        entry.lit = -entry.lit
                        assignment[abs(entry.lit)] = False
                        var_to_trail[abs(entry.lit)] = -1
                        decision.triedFalse = True

                        trailHead = decision.idx
                        break
                    } else {
                        int lit = trail.last.lit
                        assignment[abs(lit)] = Unassigned
                        var_to_trail[abs(lit)] = -1
                        erase(trail.last)
                        erase(trailStarts.last)
                    }
                    
                    if (trail.size == 0) {
                        return False
                    }

                    while (seenIdx.size > 0) {
                        int idx = seenIdx.last
                        seen[idx] = False
                        pop(seenIdx)
                    }

                    for (i = 0 up to numVars) {
                        delta[i] = False
                    }

                    continue
                }
            }

            //create learned clause and update activity scores
            int[] learned
            while (seenIdx is not empty) {
                int lit_idx = seenIdx.last
                delta[lit_idx] = True
                ...
            }

            //backjump
            ...

            //apply decay factor
            for (i = 1 up to numVars) {
                if (delta[i] == True) {
                    activity[i] = (1 - Y) + Y * activity[i]
                } else {
                    activity[i] = Y * activity[i]
                }
                delta[i] = False
            }

            continue
        }
    }
}
\end{lstlisting}
\noindent
So it's pretty straightforward,
\footnote{I struggled with debugging for half an hour because I forgot to clear out
\textit{delta} and \textit{seenIdx/seen}. I can see the headlines now: 
``Functional programmers are right about this ONE thing!!!''}
we separate our code into two branches. In the first branch where we aren't going to 
learn the clause, we just backtrack as we used to in our clause learning section, clear 
out everything related to adjusting the VSIDS and then move on, and otherwise we move on 
business as usual.

\begin{FVerbatim}
On all 1000 equations in uf20-91:

N-order:   283ms (~0.28 seconds).
AMA:       274ms (~0.27 seconds).

On all 100 equations in Flat200-479:

N-order:   85,749ms (~1.43 minutes).
AMA:       24,238ms (~24.24 seconds).

On all 1000 equations in UUF50.218.1000:

N-order:   2,313ms (~2.31 seconds).
AMA:       1,629ms (~1.63 seconds).
\end{FVerbatim}
\noindent
As you can see, doing n-order learning was terrible for our runtime. Unfortunately 
this isn't an implementation issue in our n-order learning algorithm, unless you 
consider that implementing n-order learning is an implementation issue to begin with.
You can probably guess why this is (we removed a ton of unit propagation!), and this 
is really important because it goes to show that even small initial clauses can cause 
much larger learned clauses. When I doubled $n$, the runtime was nearly the same as 
AMA, but a tad slower. 

\subsection{M-Size Relevance-Based Learning}

The kind of clause management heuristic we're going to try is m-size relevance-based 
learning, where we learn every single clause, but only delete them when they stop 
implying assignments and have more than $m$ Unassigned literals. To set this up, 
we will restart from our VSIDS code since it's not like there's anything to build off of 
from n-order learning. The paper we use uses the phrasing ``as soon as the number of
unassigned literals is greater than an integer $m$'' which sounds like we have to keep 
track of the number of unassigned literals per clause. We can do this by directly tying 
the number of Unassigned literals to the clause in a struct, which we will do. 

There is also the option of scanning the clauses in their entirety and using that information 
to choose which ones to keep and which ones to discard. Then, instead of doing this every 
conflict or whatever, we clean up the database after a certain number of conflicts.

Although adding in this struct to keep track of the number of Unassigned literals seems 
really simple, you have to consider that our watched literals implementation isn't a great 
support structure for it, since we'd need to add in a mapping from variables to clauses,
then update every single clause containing a variable, and blah blah blah, it's a lot.
There's also just the fact that honestly, I'm going to be debugging this for hours, the 
end of the quarter is draining me and I just want to graduate, and so in other words,
I'm making the executive decision to not ``waste'' more time than I have to.

So since our clause lengths are pretty small still, we'll say that once the clause 
is no longer unit, we'll get rid of it, as long as it's not part of forcing an assignment.
The implementation will be extensible, ie, it won't ``hardcode'' that the clause must be 
non-unit in order to be deleted, it will still manually count the number of Unassigned 
literals and then compare against a threshold to decide.

Here's what we need to do to get this to work: every time there's a certain number 
of conflicts, we need to scan each clause if they meet our conditions. Then for 
each ones that do, we will mark them to be deleted. Then we just delete them! So simple!
So easy! It's actually more complicated than that because of our watched literals unfortunately.
We need to remove the watch lists of each clause we're deleting, update the watch lists of all 
variables that were watching the clause we deleted, and update the watch lists of all the clauses 
that get shifted in memory, which of course, is a lot easier said than done. So here's the game 
plan: I'm gonna \textit{try} to get some sleep (it is 1:41AM), and when I wake up, maybe I'll 
have the mental fortitude to properly explain this.
\footnote{My senioritis got so bad, it evolved into full blown mental illness.}

Here's what we'll need first: we're going to implement the helper functions to help update 
our watch lists so our watched literals will work properly. So first we'll add them to our 
object:

\begin{lstlisting}
    public interface:

    private interface:
        ...
        void removeFromWatchList(int[] watchList, int idx)
        void removeClauseFromWatchLists(int clauseIdx)
        void replaceWatchIndex(int oldIdx, int newIdx)
        ...
\end{lstlisting}
\noindent
We will implement them in order.

\begin{lstlisting}
void removeFromWatchList(int[] watchList, int idx) {
    for (i = 0 up to watchList.size) {
        if (watchList[i] == idx) {
            swap(watchList[i], watchList.last)
            pop(watchList)
            break
        }
    }
}
\end{lstlisting}
This algorithm is very straightforward. Given an array of clause indices that 
we're watching, we want to remove the passed in \textit{clause index}, not 
the \textit{element} at the passed in index. We use swap-and-pop to avoid 
shifting elements, which would be a bookkeeping nightmare.

\begin{lstlisting}
void removeClauseFromWatchLists(int clauseIdx) {
    (int firstWatch, int secondWatch) = clause_to_var[clauseIdx]

    int lit1 = clauses[clauseIdx][firstWatch]
    int lit2 = clauses[clauseIdx][secondWatch]

    int varIdx1 = var_to_widx(lit1)
    int varIdx2 = var_to_widx(lit2)

    removeFromWatchList(var_to_clause[varIdx1], clauseIdx)

    if (varIdx1 != varIdx2) {
       removeFromWatchList(var_to_clause[varIdx2], clauseIdx)
    }
}
\end{lstlisting}
\noindent
This helper works by first grabbing the indices of the two literals watching 
the clause at the passed in index. We then turn them into the actual literals 
by indexing the clause. Then, we get the indices of the literals to where they 
map in \textit{var\_to\_clause}, so that way we can grab the literals' watchlists.
Finally, we remove the clause index from those literals' watchlists. This is a 
little bit confusing because of all of the indexing and grabbing the corresponding 
elements we're doing. 

\begin{lstlisting}
void replaceWatchIndex(int oldIdx, int newIdx) {
    (int firstWatch, int secondWatch) = clause_to_var(newIdx)

    int lit1 = clauses[newIdx][firstWatch]
    int lit2 = clauses[newIdx][secondWatch]

    int varIdx1 = var_to_widx(lit1)
    int varIdx2 = var_to_widx(lit2)

    int[] firstWatchList = var_to_clause[var1]
    for (i = 0 up to firstWatchList.size) {
        if (firstWatchList[i] == oldIdx) {
            firstWatchList[i] = newIdx
        }
    }

    if (varIdx1 != varIdx2) {
        int[] secondWatchList = var_to_clause[var2]
        for (i = 0 up to secondWatchList.size) {
            if (secondWatchList[i] == oldIdx) {
                secondWatchList[i] = newIdx
            }
        }
    }
}
\end{lstlisting}
\noindent
This one is pretty self explanatory. We do our indexing magic to get the watchlists 
for the clause's corresponding watch literals, and we replace an old index with 
a new index. Again, not an element \textit{at} an old index, an element \textit{equaling}
an old index.

Here is how they are used in our m-size learning algorithm:
\begin{lstlisting}
bool solveCNF() {
    //setup
    ...
    int numConflicts = 0

    while(true) {
        if (numConflicts > threshold) {
            bool[] remove

            for (i = numClauses up to clauses.size) {
                int[] clause = clauses[i]
                
                int numUnassigned = 0
                for (j = 0 up to clause.size) {
                    if (vars[abs(clause[j])] == Unassigned) {
                        numUnassigned = numUnassigned + 1
                    }
                }

                if (numUnassigned > m) {
                    remove[i - numClauses] = True
                }
            }

            for (i = 0 up to trail.size) {
                if (trail[i].reasonIdx is not optional) {
                    int reasonIdx = trail[i].reasonIdx.value

                    if (reasonIdx >= numClauses) {
                        remove[reasonIdx - numClauses] = False
                    }
                }
            }

            for (i = clauses.size down to numClauses) {
                if (remove[i - numClauses] == False) {
                    continue
                }

                int lastIdx = clauses.size - 1

                if (i == lastIdx) {
                    removeClauseFromWatchLists(i)
                    pop(clauses)
                    pop(clause_to_var)
                    continue
                }

                removeClauseFromWatchLists(i)
                swap(clauses[i], clauses[last])
                swap(clause_to_var[i], clause_to_var[last])

                replaceWatchIndex(last, idx)

                for (i = 0 up to trail.size) {
                    if (trail[i].reasonIdx is not optional) {
                        if (trail[i].reasonIdx.value == last) {
                            trail[i].reasonIdx.value = idx
                        }
                    }
                }

                pop(clauses)
                pop(clause_to_var)
            }

            numConflicts = 0
        }

        //rest of solving loop
        ...
    }
} 
\end{lstlisting}
\noindent
This is what's going on: first, we check that a certain number of conflicts have happened. 
This number is fixed. If there have been that many conflicts, we scan every clause that is 
\textbf{not} a part of our original CNF equation. If it has more than $m$ Unassigned literals, 
we mark it for deletion. In our code, our \textit{remove} array uses an indexing where we subtract 
the original number of clauses from the current clause indexing variable. This keeps the array 
smaller in our implementation. Next, we iterate through our trail, and we check if any of the 
clauses we've marked to be deleted are the reason why something was propagated. If a clause we 
want to delete is a reason clause, we can't delete it because it's still being used in our 
assignments. 

Next, we iterate from the back of the clause database up to the front, excluding the 
original clauses. If we're not supposed to remove the current clause, we move on (duh).
Otherwise, we grab the index of the last clause. If the last clause in the database is 
marked to be deleted, we can easily pop it off the database. Otherwise, we need to do 
our swap and pop nonsense. We remove the clause from its watchlists, swap the current 
clause with the last clause in the database, then we replace the watch indices. 

After that, we traverse our trail again and check if we've swapped any literal's 
reason clause index. If we have, then we need to update it. After that, we can 
remove the clause and its watched literals from the database. After this is all over,
we reset the counter for the number of conflicts and continue with the rest of the 
solving loop. 

That's a lot of work, so how does it hold up against our other implementations?
\begin{FVerbatim}
On all 1000 equations in uf20-91:

M-size:    293ms (~0.29 seconds).
N-order:   283ms (~0.28 seconds).
AMA:       274ms (~0.27 seconds).

On all 100 equations in Flat200-479:

M-size:    25,681ms (~25.68 seconds).
N-order:   85,749ms (~1.43 minutes).
AMA:       24,238ms (~24.24 seconds).

On all 1000 equations in UUF50.218.1000:

M-size:    1,937ms (~1.94 seconds).
N-order:   2,313ms (~2.31 seconds).
AMA:       1,629ms (~1.63 seconds).
\end{FVerbatim}
\noindent
We see that m-size relevance-based learning is slightly slower than AMA VSIDS, and much faster 
than n-order learning (except for uf20-91). There could be some reasons for this, such as doing 
clause deletion too frequently. Rather than doing it a fixed number of times, we could also 
perform it whenever the database grows by a certain amount, say 10\%. If we do this, our results 
speed up quite a bit on the larger equations (FLAT200.470):

\begin{FVerbatim}
On all 1000 equations in uf20-91:

M-size (10\%):  320ms (~0.32 seconds).
M-size:         293ms (~0.29 seconds).
N-order:        283ms (~0.28 seconds).
AMA:            274ms (~0.27 seconds).

On all 100 equations in Flat200-479:

M-size (10\%):  17,716ms (~17.72 seconds).
M-size:         25,681ms (~25.68 seconds).
N-order:        85,749ms (~1.43 minutes).
AMA:            24,238ms (~24.24 seconds).

On all 1000 equations in UUF50.218.1000:

M-size (10\%)   1,768ms (~1.77 seconds).
M-size:         1,937ms (~1.94 seconds).
N-order:        2,313ms (~2.31 seconds).
AMA:            1,629ms (~1.63 seconds).
\end{FVerbatim}
\noindent
Although it doesn't speed up on the smaller equations, a ~7 second speedup on larger 
equations shows that this scales better on larger SAT problems, which is great!

\subsection{Combining M-Size and K-Bounded Learning}

Due to the similarities of m-size learning and k-bounded learning, I felt like it 
didn't make much sense to implement k-bounded learning by itself, then combining it 
with m-size learning afterwards. Combining these two is really easy since all we're 
doing is checking if clauses are:
\begin{enumerate}
    \item Greater than some size $k$
    \item Have at least $m$ Unassigned literals
\end{enumerate}
Since we already do the second one in our m-size loop, it's just adding another condition
that needs to be true in order to mark a clause to be removed.

On the other hand, what isn't easy is choosing a size $k$. Not all CNF equations are 
made equal. Some equations have small clauses (like our testing sets) and learn larger 
clauses, some have large clauses and learn smaller clauses, etc. Choosing a fixed $k$ for 
all clauses is obviously a really bad idea. Make it too large and we just learn everything, 
make it too small and we forget everything. We will bring back our ``maxSize'' counter 
from n-order based learning and use that. We will use a maxSize that is twice as large as 
the largest clause in the original CNF
\footnote{Because adding a second condition to one conditional is easy, and iterating through 
the clauses to find the size of the largest clause is also easy, we will skip pseudocode.}

\begin{FVerbatim}
On all 1000 equations in uf20-91:

MK-bound:       288ms (~0.29 seconds).
M-size (10\%):  320ms (~0.32 seconds).
M-size:         293ms (~0.29 seconds).
N-order:        283ms (~0.28 seconds).
AMA:            274ms (~0.27 seconds).

On all 100 equations in Flat200-479:

MK-bound:       19,082ms (~19.08 seconds).
M-size (10\%):  17,716ms (~17.72 seconds).
M-size:         25,681ms (~25.68 seconds).
N-order:        85,749ms (~1.43 minutes).
AMA:            24,238ms (~24.24 seconds).

On all 1000 equations in UUF50.218.1000:

MK-bound:       1,740ms (~1.74 seconds).
M-size (10\%)   1,768ms (~1.77 seconds).
M-size:         1,937ms (~1.94 seconds).
N-order:        2,313ms (~2.31 seconds).
AMA:            1,629ms (~1.63 seconds).
\end{FVerbatim}
\noindent
We see that using a maxSize that's twice as large as the size of the largest original clause
causes a pretty big slowdown on Flat200-479. This indicates that we're removing clauses 
too aggressively relative to the size of the clauses we're learning.

\subsection{``Clause VSIDS''}

Now we get to the more interesting part: ``Clause VSIDS.''
\footnote{This is NOT an official or technical name as far as I am aware}
The idea is that just like in VSIDS for variables, we measure a clause's activity
and quality (LBD), and use that to determine whether or not it should be deleted. 
We can still take into account things such as length, whether or not it's being used, 
and the other factors we've been considering in other versions as well, however, the 
activity and quality are the main ones.

The paper defines this heuristic as ``the decision whether a clause should be 
deleted is based not only on the number of literals but also on its \textit{activity} in 
contributing to conflict making and on the number of decisions taken since its creation.''

In order to implement this, we'll need three additional data structures:
\begin{lstlisting}
    public interface:

    private interface:
        ...
        double[] cActivity
        int[] creationConflict
        int[] lbds
        ...
\end{lstlisting}
\noindent
\textit{cActivity} measures a learned clause's activity. We will update the activity of 
every clause involved in a conflict, if it's a learned clause. Because activity is a metric 
we use to determine whether or not a clause should be deleted, and we never delete an original 
clause, we don't need to keep track of any data for any original clauses. 
\textit{creationConflict} keeps track of the conflict number that a clause was created on.
So if we learn a clause on conflict 50, it's creation conflict is 50. We use this to measure the 
age of the clause, where $age = totalConflicts - creationConflict$.
Lastly, \textit{lbds} keeps track of a clause's LBD upon learn time. Even though a clause's LBD 
changes as our decision trail changes, we want to keep the original LBD from when the clause was 
learned. This is because we want to know what decision levels led to the conflict in the first place, 
and recomputing it describes the how ``useful'' the learned clause is under the current assignment, 
not when it was actually created. Some SAT solvers do use dynamic LBD computations, but it's a 
different heuristic.

In order to keep track of these metrics for each learned clause, we need to update our 
VSIDS code, as well as some other parts of our conflict resolution algorithm. Luckily, 
these changes are very minimal:
\begin{lstlisting}
bool solveCNF() {
    //setup
    ...
    int totalConflicts = 0
    int curSize = numClauses

    double activityIncrease = 1.0
    double activityGrowth = 1.05
    double maxIncrease = 1e100
    double resetIncrease = 1e-100

    while (true) {
        if (activityIncrease > maxIncrease) {
            for each (clauseActivity in cActivity) {
                clauseActivity *= resetIncrease
            }
            activityIncrease *= resetIncrease
        }

        //Current clause deletion
        if (clauseDatabase.size grew by X%) {
            ...

            for (i = clauses.size down to numClauses) {
                if (remove[i - numClauses] == False) {
                    continue
                }

                int lastIdx = clauses.size - 1

                if (i == lastIdx) {
                    removeClauseFromWatchLists(i)
                    pop(clauses)
                    pop(clause_to_var)

                    pop(cActivity)
                    pop(creationConflict)
                    pop(lbds)

                    continue
                }

                removeClauseFromWatchLists(i)
                swap(clauses[i], clauses[last])
                swap(clause_to_var[i], clause_to_var[last])

                swap(cActivity[i], cActivity[last])
                swap(creationConflict[i], creationConflict[last])
                swap(lbds[i], lbds[last])

                ...

                pop(clauses)
                pop(clause_to_var)

                pop(cActivity)
                pop(creationConflict)
                pop(lbds)
            }

            curSize = clauses.size
        }

        //assignment
        ...

        if (contradiction occurred) {
            totalConflicts = totalConflicts + 1
            ...

            for (int trailIdx = trail.size down to 0) {
                ...
                int reasonIdx = trail[trailIdx].reasonIdx.value;
                if (reasonIdx >= numClauses) {
                    cActivity[reasonIdx - numClauses] += activityIncrease
                }
            }
            //rest of VSIDS and LBD calculation
            ...

            //backjump and learn 
            ... 

            append(cActivity, 0.0)
            append(creationConflict, totalConflicts)
            append(lbds, lbd)
            activityIncrease *= activityGrowth
            ...
        }
    }
}
\end{lstlisting}
\noindent
In our clause deletion code, we just added code to update our new data structures
to be consistent with the rest of the clause database data structures. And then 
in the VSIDS calculation, we added code to update a clause's activity if it's 
a learned clause, and at the end of the loop, we just add extra entries into each 
of the data structures so they grow with the database. 

We use an activity increase and an activity growth variable to simulate the effects 
of an activity decay and bump. The idea is that instead of bumping only the clauses 
that were involved in a conflict and decaying the rest, we bump only the clauses 
involved in a conflict in a way that is as if we decayed the others. By increasing 
the activity increase by some growth amount every conflict, when we add to a clause's 
activity, we add an ever increasing amount, so that way, more recent clauses get 
larger and larger activities and less recent clauses don't have such huge increases. 
The only caveat to this is that because conflicts happen all the time, we need to 
reset the activity increase once it grows to a stupidly large size. This can be thought 
of as our decay step.

The next step is to now use these metrics in our clause deletion code. We could do 
a bunch of experiments of how each metric impacts runtime and clause deletion, but 
frankly, I don't want to. Instead, here's what workflow we'll use:
\begin{enumerate}
    \item Find which clauses that can be deleted by LBD, age, and activity
    \item Find which clauses can be deleted via m+k-bounded learning
    \item Protect clauses with an LBD \verb|<=| 2
    \item Protect clauses that are currently reason clauses
    \item Delete the ``bad'' ones
\end{enumerate}
\noindent
To determine what a ``bad'' clause means, we need to think of what the goals of 
clause deletion is. For instance, do we want to delete a fixed number of clauses, 
or should it be more dynamic, as in, there isn't a quota to fulfill, instead we look at 
each clause and go ``you suck, get out'' or ``you're good, stay?'' This idea is pretty 
attractive, but there is something we should consider: what if there's not enough bad 
clauses relative to our database size, or what if there's too many bad clauses relative 
to our database size? Ie, what if we mark too little clauses for deletion and our database 
is still large, as well as marking too many clauses for deletion and then our database 
becomes super small again? To be honest, if there aren't enough ``bad'' clauses, to delete,
to me it seems like most of the database is really good, so I wouldn't mark \textit{extra}
clauses for deletion, however, capping how many clauses we delete does seem like a good idea.

Because we have a clear marker of when clause deletion should occur (when the database 
grows by 10\%), we will say that the number of clauses being deleted can't exceed this 
10\% threshold. So if our original database has 100 clauses and it grows to 110 clauses,
we can't delete more than 9 clauses. Supposing we delete 0, and the database grows to 
121, we can't delete more than 10 clauses, and so on. Another advantage of this is 
that the maximum number of clauses we can delete grows/shrinks with our threshold, so 
if in a later experiment we found that it's more beneficial to run the deletion algorithm 
every time the database grows by 20\%, we don't have to manually update the number of 
clauses we're allowed to delete as well.

In order to decide which clauses to delete, we want to decide a threshold for our 
metrics. I used the ones in the bottom 25th percentile for their respetive metric, 
so high LBD, high age, and low activity. We then select the clauses in those 
bottom metrics, and if they have at least 2 of these negative attributes, we mark 
them to be deleted. 

After that, we do our regular m+k bounded learning, then afterwards, we mark clauses 
that have low LBD (\verb|<=| 2), and clauses that are reason clauses as protected, 
so they can't be deleted even if they had at least two bad attributes.

The code is as follows:
\begin{lstlisting}
bool solveCNF() {
    //setup
    ... 

    while (true) {
        //Current clause deletion
        if (clauseDatabase.size grew by X%) {
            int numCanDelete = clauseDatabase.size - curSize

            int[] candidateIdx
            for (i = 0 up to clauseDatabase.size - numClauses) {
                candidateIdx[i] = i
            }

            int cutoff = candidateIdx.size * 0.25

            sortIncreasing(cActivity)

            double activityCutoff = cActivity[cutoff]

            sortIncreasing(conflictCreation)

            int ageCutoff = conflictCreation[cutoff]

            sortDecreasing(lbds)

            int lbdCutoff = lbds[cutoff]

            for (i = numClauses up to clauses.size) {
                bool lowActivity = cActivity[i] <= activityCutoff
                bool highAge = creationConflict[i] <= ageCutoff
                bool highlbd = lbds[i] >= lbdCutoff

                remove[i - numClauses] = lowActivity + highAge + highlbd >= 2
            }

            //rest of the database deletion
            ...
        }

        //rest of solving
        ...
    }

}
\end{lstlisting}

\begin{FVerbatim}
On all 1000 equations in uf20-91:

Clause VSIDS:   347ms (~0.35 seconds).
MK-bound:       288ms (~0.29 seconds).
M-size (10\%):  320ms (~0.32 seconds).
M-size:         293ms (~0.29 seconds).
N-order:        283ms (~0.28 seconds).
AMA:            274ms (~0.27 seconds).

On all 100 equations in Flat200-479:

Clause VSIDS:   19,181ms (~19.18 seconds).
MK-bound:       19,082ms (~19.08 seconds).
M-size (10\%):  17,716ms (~17.72 seconds).
M-size:         25,681ms (~25.68 seconds).
N-order:        85,749ms (~1.43 minutes).
AMA:            24,238ms (~24.24 seconds).

On all 1000 equations in UUF50.218.1000:

Clause VSIDS:   1,804ms (~1.80 seconds).
MK-bound:       1,740ms (~1.74 seconds).
M-size (10\%)   1,768ms (~1.77 seconds).
M-size:         1,937ms (~1.94 seconds).
N-order:        2,313ms (~2.31 seconds).
AMA:            1,629ms (~1.63 seconds).
\end{FVerbatim}
\noindent
To be honest, I'm a little bit disappointed in these results. I was really hopeful 
that this heuristic would have yielded more gains in speed, not runtime. Was my 
implementation perfect? No. Could that have contributed to the increased runtime?
Yes, but I have a feeling that it's also the testing sets being used. Again, we 
can't necessarily draw a conclusion on \textit{all} SAT problems based on this 
data, but we can draw the conclusion that the solve is worse for \textit{these}
SAT problems.

There were some things I tried to see what could be causing this. I tried adjusting 
how aggressive/loose the badness score had to be in order to be marked to be deleted. 
A lower badness score threshold correlated with a higher runtime. Then I tried removing 
the cap on deletion altogether, and that didn't impact the runtime very much. Then I tried 
adjusting the percentile to be the 20th, 25th, and 50th, and they all had very similar runtimes 
as well. I even tried messing with what we consider a bad clause, and that impacted 
the runtime, in a bad way. Lastly, I tried removing mk-bound learning, and that also didn't 
make much of an impact. I then removed the k-bounded learning afterwards, and that also didn't 
seem to help. No matter what I tried, it appears that this heuristic was hurting the runtime 
on all accounts. Perhaps something to be explored at a later date.

\subsection{Clause Database Management - Results}

Here are the results for all the kinds of clause database management compared to 
the last implementation without any kind of management:

\begin{FVerbatim}
On all 1000 equations in uf20-91:

Clause VSIDS:               347ms (~0.35 seconds).
MK-bound:                   288ms (~0.29 seconds).
M-size (10\%):              320ms (~0.32 seconds).
M-size:                     293ms (~0.29 seconds).
N-order:                    283ms (~0.28 seconds).

AMA:                        274ms (~0.27 seconds).
EMA (X = 1, Y = 0.95):      266ms (~0.27 seconds).

Improved CDCL:              284ms (~0.28 seconds).
First UIPs:                 406ms (~0.41 seconds).
Backjumping:                380ms (~0.38 seconds).
Clause Learning:            425ms (~0.43 seconds).
Iterative DPLL (AbP):       328ms (~0.33 seconds).
Iterative DPLL (PbA):       324ms (~0.32 seconds).
Watched Literals:           349ms (~0.35 seconds).
Pure Literal Elimination:   1548ms (~1.55 seconds).
Unit Propagation:           868ms (~0.87 seconds).
Recursive Backtracking:     1,540,613ms (~25.68 minutes).
Brute Force:                1,770,851ms (~29.51 minutes).
Short-Circuited (Basic):    49,874ms (~50 seconds).
Parallel:                   10,589ms (~10.5 seconds).

On all 100 equations in Flat200-479:

Clause VSIDS:   19,181ms (~19.18 seconds).
M-size (10\%):  17,716ms (~17.72 seconds).
AMA:            24,238ms (~24.24 seconds).

On all 1000 equations in UUF50.218.1000:

Clause VSIDS:               1,804ms (~1.80 seconds).
MK-bound:                   1,740ms (~1.74 seconds).
M-size (10\%)               1,768ms (~1.77 seconds).
M-size:                     1,937ms (~1.94 seconds).
N-order:                    2,313ms (~2.31 seconds).

AMA:                        1,629ms (~1.63 seconds).
EMA (X = 1, Y = 0.95):      1,504ms (~1.50 seconds).

Improved CDCL:              2,385ms (~2.39 seconds).
First UIPs:                 14,012ms (~14.01 seconds).
Backjumping:                11,291ms (~11.29 seconds).
Clause Learning:            9,652ms (~9.65 seconds).
Iterative DPLL (AbP):       6,369ms (~6.37 seconds).
Iterative DPLL (PbA):       6,365ms (~6.37 seconds).
Watched Literals:           7,130ms (~7.13 seconds).
Pure Literal Elimination:   231,107ms (~3.85 minutes).
Unit Propagation:           140,041ms (~2.34 minutes).
\end{FVerbatim}
\noindent

Because of how well M-size (10\%) performed, I will keep it in the overall results 
when compared with every other solver we've created so far:

\begin{FVerbatim}
On all 1000 equations in uf20-91:

Clause VSIDS:               347ms (~0.35 seconds).
M-size (10\%):              320ms (~0.32 seconds).

AMA (Y = 0.75 or 0.99):     274ms (~0.27 seconds).
EMA (X = 1, Y = 0.95):      266ms (~0.27 seconds).

Improved CDCL:              284ms (~0.28 seconds).
First UIPs:                 406ms (~0.41 seconds).
Backjumping:                380ms (~0.38 seconds).
Clause Learning:            425ms (~0.43 seconds).
Iterative DPLL (AbP):       328ms (~0.33 seconds).
Iterative DPLL (PbA):       324ms (~0.32 seconds).
Watched Literals:           349ms (~0.35 seconds).
Pure Literal Elimination:   1548ms (~1.55 seconds).
Unit Propagation:           868ms (~0.87 seconds).
Recursive Backtracking:     1,540,613ms (~25.68 minutes).
Brute Force:                1,770,851ms (~29.51 minutes).
Short-Circuited (Basic):    49,874ms (~50 seconds).
Parallel:                   10,589ms (~10.5 seconds).

On all 100 equations in Flat200-479:

Clause VSIDS:               19,181ms (~19.18 seconds).
M-size (10\%):              17,716ms (~17.72 seconds).
AMA:                        24,238ms (~24.24 seconds).

On all 1000 equations in UUF50.218.1000:

Clause VSIDS:               1,804ms (~1.80 seconds).
M-size (10\%)               1,768ms (~1.77 seconds).

AMA:                        1,629ms (~1.63 seconds).
EMA (X = 1, Y = 0.95):      1,504ms (~1.50 seconds).

Improved CDCL:              2,385ms (~2.39 seconds).
First UIPs:                 14,012ms (~14.01 seconds).
Backjumping:                11,291ms (~11.29 seconds).
Clause Learning:            9,652ms (~9.65 seconds).
Iterative DPLL (AbP):       6,369ms (~6.37 seconds).
Iterative DPLL (PbA):       6,365ms (~6.37 seconds).
Watched Literals:           7,130ms (~7.13 seconds).
Pure Literal Elimination:   231,107ms (~3.85 minutes).
Unit Propagation:           140,041ms (~2.34 minutes).
\end{FVerbatim}
\noindent
Sorry that this section is really bloated.